% dynamics/realization_aware_dynamics-DE.tex
\chapter{Realisierungsbewusste Vorhersage und Vergleich}
\label{chap:realization-aware-dynamics}

\section*{Welche Geometrie geht in das Modell ein?}

Beabsichtigte Alignment-Geometrie, ihre Metric Realization, gebautes Asset,
physischer Zustand und beobachtete Zustandsschätzung sind verbunden, aber
nicht austauschbar. Eine Koordinatentransformation ändert eine
Repräsentation; sie baut kein Asset. Eine Vermessungspolylinie ist eine
beobachtungstragende Repräsentation und keine physische Wahrheit.

\begin{principle}[Prinzip realisierungsbewusster Dynamik]
\label{princ:realization-aware-dynamics}
Eine Vorhersage benennt, welche qualifizierte Realisierung oder beobachtete
Zustandsschätzung in das Antwortmodell \(M\) eingeht. Das Ergebnis bleibt
eine Vorhersage von \(M\), keine Beobachtung.
\end{principle}

\section{Typisierte Vergleichskette}

\begin{definition}[Realisierter Laufzeitzustand]
\label{def:realized-runtime-state}
Ein realisierter Laufzeitzustand ist hier eine Eingabe des Antwortmodells, die
aus einer ausdrücklich identifizierten Realisierung oder beobachteten
Zustandsschätzung abgeleitet und durch Objektidentität, Zeit, Frame, Methode,
Provenienz, Anwendbarkeit und Unsicherheit qualifiziert ist. Er ist keine
neue kanonische Zustandsklasse.
\end{definition}

Für die Quelle \(q\in\{\mathrm{Absicht},\mathrm{Realisierung},
\mathrm{Beobachtung}\}\) gilt
\[
\widehat S_q
\xrightarrow[\;Q_M,I_M\;]{\mathcal D_M}
x_{M,q}^{\mathrm{pred}}(t).
\]
Nur gleichartige Ausgaben dürfen verglichen werden. Eine gemessene Antwort
\(x^{\mathrm{obs}}(t)\) bleibt getrennt qualifiziert:
\[
\Delta_M(t)=
\mathcal C\!\left(x_{M,q}^{\mathrm{pred}}(t),
x^{\mathrm{obs}}(t)\mid
\text{gleiches Subjekt, gleiche Zeit, gleicher Frame und Größe}\right).
\]
\(\Delta_M\) ist ein Vergleichsergebnis, keine Ursache oder Entscheidung.

\section{Realisierungsfilterung}

\begin{proposition}[Prinzip des Realisierungsfilters]
\label{prop:realization-filter-principle}
Eine Änderung der qualifizierten Modelleingabe kann die Vorhersage ändern und
zugleich die Alignment Identity bewahren. Die Änderung darf erst dann
Konstruktion, Degradation oder Messung zugeschrieben werden, wenn Evidenz
diese Erklärungen unterscheidet.
\end{proposition}

Bei der Abweichung von CAD-Kurve und Vermessungspolylinie dürfen
CAD-Realisierung und polylinienabgeleitete Schätzung jeweils dasselbe \(M\)
speisen. Die Differenz der Ausgaben ist ein Sensitivitätsergebnis. Eine
kausale Zuordnung benötigt weitere Evidenz.

\section{Kalibrierungs- und Unsicherheitsgrenze}

Ein Modellparameter gehört zu \(Q_M\). Kalibrierungsevidenz stützt Auswahl
oder Schätzung eines Parameters für einen angegebenen Bereich.
Beobachtungsunsicherheit qualifiziert Messevidenz; Konstruktionsvariabilität
betrifft das Asset; operative Variabilität die Nutzung; ein numerisches
Residuum eine Rechnung; Modellformunsicherheit die Angemessenheit von \(M\).
Diese Quellen dürfen nicht allein deshalb zu einem Fehlerterm verschmolzen
werden, weil sie alle einen Vergleich beeinflussen.

Die Übereinstimmung mit einer Beobachtung kann eine lokale
Kalibrierungsaussage stützen. Sie begründet keine universelle Modellgültigkeit.

\begin{definition}[Realisierungsbewusste Zulässigkeit]
\label{def:realization-aware-admissibility}
Realisierungsbewusste Zulässigkeit ist ein Bewertungsergebnis aus einer
identifizierten Regel und einem Zweck auf qualifizierter Vergleichsevidenz.
Sie ist keine intrinsische Eigenschaft, die \(\mathcal D_M\) erzeugt.
\end{definition}

\begin{definition}[Laufzeithülle]
\label{def:runtime-envelope}
Eine Laufzeithülle ist eine erklärte Grenzwertmenge aus einer identifizierten
Quelle, mit Einheiten, Anwendbarkeit und Gültigkeit. Das Zustandsmodell selbst
liefert sie nicht.
\end{definition}

Damit entstehen Eingaben für eine Optimierung, aber noch kein Ziel. Die
Verantwortung dafür wird nun ausdrücklich gemacht.
